\section{Introduction}
\begin{frame}[fragile]{Stochastic Multiplayer Games (SMGs)}
\begin{columns}
    \column{.5\textwidth}
        \tikzset{
          myNode/.style={draw=black, text=black, minimum size=1cm},
          highlightLine/.style={draw=black},
          highlightNode/.style={fill=black},
          singleLetterLabel/.style={color=black},
          fractionLabel/.style={color=black}
        }
        %--- Slide 1: Node text blue
        \only<1>{
          \tikzset{
            myNode/.append style={text=orange, draw=black} % text=blue, outline=black
          }
        }
        %--- Slide 2: Node outline blue
        \only<2>{
          \tikzset{
            myNode/.append style={text=black, draw=orange} % text=black, outline=blue
          }
        }
        %--- Slide 3: Single-letter labels blue
        \only<3>{
          \tikzset{
            singleLetterLabel/.append style={color=orange}
          }
        }
        %--- Slide 4: Fraction labels, lines, small circles blue
        \only<4>{
          \tikzset{
            fractionLabel/.append style={color=orange},
            highlightLine/.append style={draw=orange},
            highlightNode/.append style={fill=orange}
          }
        }
        \begin{tikzpicture}[>=stealth, 
                            every node/.style={font=\small}, 
                            node distance=3cm]
          \node[myNode, shape=rectangle] (S1)                  {$s_1$};
          \node[myNode, shape=circle,    right=5cm of S1] (S2) {$s_2$};
          \node[myNode, shape=circle, 
                below=3cm of $(S1)!0.5!(S2)$] (S3)             {$s_3$};
        
          \draw[highlightLine] (S1) 
            -- node[singleLetterLabel, below left, pos=0.5] {$a$} 
               ($(S1)!0.5!(S3)$);
          \filldraw[highlightNode] ($(S1)!0.5!(S3)$) circle (2pt);
          \draw[->, highlightLine] 
            ($(S1)!0.5!(S3)$) 
            -- node[fractionLabel, below left, pos=0.5] {$\frac{2}{3}$} 
               (S3);
          \draw[->, highlightLine]
            ($(S1)!0.5!(S3)$)
            -- node[fractionLabel, below, pos=0.75] {$\frac{1}{3}$}
               (S2);
          \draw[highlightLine] (S2) 
            -- node[singleLetterLabel, below right, pos=0.5] {$a$} 
               ($(S2)!0.5!(S3)$);
          \filldraw[highlightNode] ($(S2)!0.5!(S3)$) circle (2pt);
          \draw[->, highlightLine]
            ($(S2)!0.5!(S3)$)
            -- node[fractionLabel, below right, pos=0.5] {$\frac{1}{2}$}
               (S3);
          \draw[->, highlightLine]
            ($(S2)!0.5!(S3)$)
            -- node[fractionLabel, below, pos=0.75] {$\frac{1}{2}$}
               (S1);
          \draw[->, highlightLine]
            (S1) to[bend left=30]
            node[singleLetterLabel, above, pos=0.25] {$b$}
            node[fractionLabel,    above, pos=0.75] {$1$}
            coordinate[pos=0.5] (S1S2)
            (S2);
          \filldraw[highlightNode] (S1S2) circle (2pt);
          \draw[->, highlightLine]
            (S2) to
            node[singleLetterLabel, above, pos=0.25] {$b$}
            node[fractionLabel,    above, pos=0.75] {$1$}
            coordinate[pos=0.5] (S2S1)
            (S1);
          \filldraw[highlightNode] (S2S1) circle (2pt);
          \draw[->, highlightLine] 
            (S3) edge[loop below]
            coordinate[pos=0.5] (S3S3)
            node[singleLetterLabel, right, pos=0.2] {$c$}
            node[fractionLabel,     left,  pos=0.8] {$1$} 
            (S3);
          \filldraw[highlightNode] (S3S3) circle (2pt);
        \end{tikzpicture}

    \column{0.5\textwidth}
        \vspace{4mm}
        
        A SMG $\mathcal{G}$ consists of:
        \begin{itemize}
            \item<1-> States $S=\set{s_1, s_2, s_3}$,\\associated with
            \item<2-> Players $\Pi=\set{\text{square}, \text{circle}}$\\
                (or $\set{0, 1}$)
            \item<3-> Actions $A=\set{a, b}$
            \item<4-> A transition function $\Delta$
        \end{itemize}
        \vspace{3mm}
        \uncover<5->{
        A \textit{strategy} is a function that tells a player what action to choose for every possible path.
        }
        
        \vspace{2mm}
        \uncover<6->{
        Strategies allow us to reason about the probability of certain outcomes.
        }
\end{columns}
\end{frame}




\begin{frame}{Positive Reachability}
\begin{columns}
    \column{.5\textwidth}
\begin{tikzpicture}[>=stealth, node distance=3cm, every node/.style={font=\small}]

% Nodes
\node[draw, minimum size=1cm, shape=rectangle] (S1) {$s_1$};
\node[draw, minimum size=1cm, shape=circle,right= 5cm of S1] (S2) {$s_2$};
\node[draw, minimum size=1cm, shape=circle, below= 3cm of $(S1)!0.5!(S2)$] (S3) {$s_3$};

\draw[-] (S1) -- node[below left, pos=0.5] {$a$} ($(S1)!0.5!(S3)$);
\filldraw[black] ($(S1)!0.5!(S3)$) circle (2pt); % Dot at split point
\draw[->] ($(S1)!0.5!(S3)$) -- node[below left, pos=0.5] {$\frac{2}{3}$} (S3);
\draw[->] ($(S1)!0.5!(S3)$) -> node[below, pos=0.75] {$\frac{1}{3}$} (S2);
\draw[-] (S2) -- node[below right, pos=0.5] {$a$} ($(S2)!0.5!(S3)$);
\filldraw[black] ($(S2)!0.5!(S3)$) circle (2pt); % Dot at split point
\draw[->] ($(S2)!0.5!(S3)$) -- node[below right, pos=0.5] {$\frac{1}{2}$} (S3);
\draw[->] ($(S2)!0.5!(S3)$) -- node[below, pos=0.75] {$\frac{1}{2}$} (S1);

\draw[->] (S1) to[bend left=30] node[above, pos=0.25] {$b$} 
node[above, pos=0.75] {$1$} coordinate[pos=0.5] (S1S2) (S2);
\filldraw[black] (S1S2) circle (2pt); % Dot halfway along the curve

\draw[->] (S2) -> node[above, pos=0.25] {$b$} node[above, pos=0.75] {$1$} coordinate[pos=0.5] (S2S1) (S1);
\filldraw[black] (S2S1) circle (2pt); % Dot halfway along the curve

\draw[->] (S3) edge[loop below]
coordinate[pos=0.5] (S3S3)
node[right, pos=0.2] {$c$}
node[left, pos=0.8] {$1$} (S3);
\filldraw[black] (S3S3) circle (2pt); % Dot halfway along the curve

\end{tikzpicture}
\column{.5\textwidth}
\large
\vspace{5mm}

\uncover<+->{
Starting at $s_1$, does player $0$ have a strategy of reaching the state $s_3$ eventually with a non-zero probability?
}
\vspace{5mm}

\uncover<+->{
$\Rightarrow$ Yes, pick $a$ as the first action.\\
}
\vspace{5mm}

\normalsize
\uncover<+->{
In general:\\
Determine the set $\Reach_i(T, \mathcal{G})$.
This is called \alert{positive reachability}.
}
\vspace{5mm}

\uncover<+->{
In this case:
\\$\Reach_0(\set{s_3}, \mathcal{G})=\set{s_1, s_2, s_3}$
}
\end{columns}
\end{frame}

\begin{frame}{Almost-Sure Reachability}
\begin{columns}
    \column{.5\textwidth}
\begin{tikzpicture}[>=stealth, node distance=3cm, every node/.style={font=\small}]

% Nodes
\node[draw, minimum size=1cm, shape=rectangle] (S1) {$s_1$};
\node[draw, minimum size=1cm, shape=circle,right= 5cm of S1] (S2) {$s_2$};
\node[draw, minimum size=1cm, shape=circle, below= 3cm of $(S1)!0.5!(S2)$] (S3) {$s_3$};

\draw[-] (S1) -- node[below left, pos=0.5] {$a$} ($(S1)!0.5!(S3)$);
\filldraw[black] ($(S1)!0.5!(S3)$) circle (2pt); % Dot at split point
\draw[->] ($(S1)!0.5!(S3)$) -- node[below left, pos=0.5] {$\frac{2}{3}$} (S3);
\draw[->] ($(S1)!0.5!(S3)$) -> node[below, pos=0.75] {$\frac{1}{3}$} (S2);
\draw[-] (S2) -- node[below right, pos=0.5] {$a$} ($(S2)!0.5!(S3)$);
\filldraw[black] ($(S2)!0.5!(S3)$) circle (2pt); % Dot at split point
\draw[->] ($(S2)!0.5!(S3)$) -- node[below right, pos=0.5] {$\frac{1}{2}$} (S3);
\draw[->] ($(S2)!0.5!(S3)$) -- node[below, pos=0.75] {$\frac{1}{2}$} (S1);

\draw[->] (S1) to[bend left=30] node[above, pos=0.25] {$b$} 
node[above, pos=0.75] {$1$} coordinate[pos=0.5] (S1S2) (S2);
\filldraw[black] (S1S2) circle (2pt); % Dot halfway along the curve

\draw[->] (S2) -> node[above, pos=0.25] {$b$} node[above, pos=0.75] {$1$} coordinate[pos=0.5] (S2S1) (S1);
\filldraw[black] (S2S1) circle (2pt); % Dot halfway along the curve

\draw[->] (S3) edge[loop below]
coordinate[pos=0.5] (S3S3)
node[right, pos=0.2] {$c$}
node[left, pos=0.8] {$1$} (S3);
\filldraw[black] (S3S3) circle (2pt); % Dot halfway along the curve

\end{tikzpicture}
\column{.5\textwidth}
\large
\vspace{5mm}

\uncover<+->{
Starting at $s_1$, does player $0$ have a strategy of reaching the state $s_3$ eventually with probability $1$?
}
\vspace{5mm}

\uncover<+->{
$\Rightarrow$ Yes, always pick $a$.
}
\vspace{5mm}

\uncover<+->{
\normalsize
In general:\\
Determine the set $\ASReach_i(T, \mathcal{G})$.
This is called \alert{almost-sure reachability}.
\vspace{5mm}
}

\uncover<+->{
In this case:
\\$\ASReach_0(\set{s_3}, \mathcal{G})=\set{s_1, s_2, s_3}$
}
\end{columns}
\end{frame}

\begin{frame}{Contributions}
\begin{itemize}[<+->]
    \item Detailed formalization of an algorithm based on Strong- and Weak-Attractors by Hahn et al.
    \item Proof of correctness
    \item Performant implementation of the algorithm in Storm
    \item Experimental evaluation
\end{itemize}
    
\end{frame}